\originalchapter{白矮星冷却与恒星风习题}
本章完整对应8.901 PS6.1--5，均为独立补写。原作业中经验关系按其年代和单位使用；额外的包层模型、标准质量和初始质量分别声明，不把模型演算冒充当前精密观测。

\src{8.901 PS6.1，40 Eri B白矮星，(a)--(d)。}
\begin{exercise}
40 Eri B和C绕行周期247.9 yr、视差0.201角秒、相对轨道半长轴6.89角秒、$a_B/a_C=0.37$；A约在400 AU外，本题只处理B+C。
(a)求B、C质量；(b)B的绝对热星等9.6，求光度；
(c)$T_{\rm eff}=16900$ K，求半径、平均密度；
(d)与Sirius B的$R=0.008\Rs$、$\bar\rho=3.0\times10^6\ \mathrm{g\,cm^{-3}}$比较并说明密度差。
\end{exercise}
\begin{solution}
距离$D=1/0.201=4.975$ pc，半长轴$a=6.89D=34.279$ AU。在太阳质量、AU、yr单位中$M_B+M_C=a^3/P^2=0.6554\Ms$。质心关系$M_Ba_B=M_Ca_C$给$M_C/M_B=0.37$，所以$M_B=0.4784\Ms,M_C=0.1770\Ms$；位置比不能误当同方向质量比。

明确取热星等零点$M_{{\rm bol},\odot}=4.74$，$L_B/\Ls=10^{-0.4(9.6-4.74)}=0.01138$。有效温度定义给$R_B/\Rs=\sqrt{L_B/\Ls}(5772/16900)^2=0.01244$，平均密度$3M_B/(4\pi R_B^3)=3.50\times10^5\ \mathrm{g\,cm^{-3}}$；平均密度不是中心密度。

Sirius B平均密度约高8.6倍、半径更小。非相对论电子简并支持给大致$R\propto M^{-1/3}$，较大质量可有较小半径和更高密度；组成、温度及相对论修正亦影响。不能从“更亮”直接推“更密”。这里保持题给2006年轨道参数，不作为40 Eri当前更新质量。
\end{solution}

\src{8.901 PS6.2，结晶与潜热冷却，(a)--(d)。}
\begin{exercise}
纯碳白矮星$M=0.6\Ms,R=0.012\Rs$。相邻离子尺度定义$4\pi r_i^3/3=Am_p/\rho$，Coulomb参数$\Gamma=Z_i^2e^2/(r_i\kb T)$。
(a)求$r_i$；(b)以$\Gamma=160$估内部结晶温度；
(c)估此温度的光度，包层质量分数$X=0,Y=0.9,Z=0.1$；
(d)每核释放$\kb T$潜热时，可维持此光度多久？
\end{exercise}
\begin{solution}
(a) 平均密度$\bar\rho=4.895\times10^5\ \mathrm{g\,cm^{-3}}$，碳$A=12,Z_i=6$，得$r_i=2.139\times10^{-10}$ cm。
(b) Gaussian单位$e^2=\alpha\hbar c$，$T_c=36e^2/(160r_i\kb)=1.758\times10^6$ K。实际核心密度大于平均密度，结晶会分区发生；此单温度是题设平均模型估计。

(c) 题面没有给包层不透明度、边界或$L(T_c)$；成分本身不唯一决定光度。可明确补入Mestel型示例：理想包层$P=\rho\kb T/(\mu m_p)$，$\mu^{-1}=3Y/4+Z/2$，辐射Kramers不透明度$\kappa=\kappa_0\rho T^{-7/2}$，静力质量近$M$，表面$P,T$近零。将扩散式除以静力式，有
\[
 \frac{\dd T}{\dd P}=\frac{3\kappa_0L\mu m_p}{16\pi a_{\rm rad}cGM\kb}
 P T^{-15/2}.
\]
积分得
\[
 P^2=\frac{64\pi a_{\rm rad}cGM\kb}{51\kappa_0L\mu m_p}T^{17/2}.
\]
以理想气体压力等于非相对论电子压力$K_{\rm NR}\rho^{5/3}$作简化匹配，$\rho_b=[\kb T_c/(\mu m_pK_{\rm NR})]^{3/2}$，从而
\[
 L=\frac{64\pi a_{\rm rad}cGM}{51\kappa_0}
 \left(\frac{\mu m_p}{\kb}\right)^4K_{\rm NR}^3T_c^{7/2}.
\]
示例额外定义$\kappa_0=4.6712\times10^{23}$（相应cgs单位），$\mu=1.3793,\mu_e=2$，实际得到$L=6.81\times10^{-7}\Ls$。这不是题面给的系数；$L\propto\kappa_0^{-1}$，更透明包层、不同匹配或实际核密度都会显著改值。通式明确表达缺少哪些物理输入，不将示例说成原题唯一光度。

(d) 核数$N_i=M/(12m_p)$，潜热$E_{\rm lat}=N_i\kb T_c$，固定光度时间$\tau=E_{\rm lat}/L$。上述示例给$1.75\times10^{11}$ yr，巨大值提示极简不透明包层与平均密度组合不可作为真实冷却年龄：光度不会恒定，核心结晶结构不同，还存在普通热容量等能量。原题条件下可给$r_i,T_c$与条件化$L,\tau$；若需真实冷却预测必须补实际包层输运和结晶过程，不能复制未取得的教材冷却图补空。
\end{solution}

\src{8.901 PS6.3，白矮星计时与银河年龄。}
\begin{exercise}
本地白矮星光度函数截断$L\simeq10^{-4.4}\Ls$。初始主序质量$0.6<M/\Ms<10$的形成白矮星前时间满足
\[
 \log_{10}\tau_{\rm pn}(\mathrm{yr})=9.921-3.6648\ell+1.9697\ell^2-0.9369\ell^3,
 \quad\ell=\log_{10}(M/\Ms).
\]
随后冷却式为$\tau_{\rm wd}=8.8\times10^6(A/12)^{-1}(M_{\rm wd}/\Ms)^{5/7}(\mu/2)^{-2/7}(L/\Ls)^{-5/7}$ yr。按原题纯碳$A=\mu=12$及标准质量，估银河盘年龄；假设样本足够大，最暗端来自能成为白矮星的最大初始质量，并解释原因。晕的截断$T_{\rm eff}=4000$ K，取标准质量、半径再估年龄。
\end{exercise}
\begin{solution}
总年龄是前身演化$\tau_{\rm pn}$加冷却$\tau_{\rm wd}$。较高质量前身较早变成白矮星，有更长的冷却时间；足够大样本才可能采到这种暗端物体。最大前身质量未给，示例取$8\Ms$，代多项式得$3.37\times10^7$ yr；白矮星另取$0.6\Ms$，不能继续用$8\Ms$算冷却。

原题写$A=\mu=12$；这里忠实按这一给定冷却参数，盘冷却$5.088\times10^9$ yr，加前身时间约5.12 Gyr。晕若取$R=0.012\Rs$，则$L/\Ls=(0.012)^2(4000/5772)^4=3.321\times10^{-5}$，冷却约5.791 Gyr，总约5.82 Gyr。

原$\mu$称平均分子量，但充分电离纯碳气体的通常粒子分子量是12/7，不是12。若把冷却式$\mu$解释为这种气体分子量，时间比原计算大$7^{2/7}\simeq1.74$；不能悄悄混用。真实银河年龄还受初始—最终质量关系、形成历史、包层、结晶潜热和样本完备性影响。这里是给定时钟的条件化结果；题面已给截断事实，无需复制未取得的教材光度函数图。
\end{solution}

\src{8.901 PS6.4，热亮星的风，(a)--(b)。}
\begin{exercise}
题给$\log_{10}(\dot M v_\infty R^{1/2})=-1.37+2.07\log_{10}(L/10^6)$，$M,R,L$为太阳单位，$\dot M$为$\Ms$/yr，$v_\infty$为km/s。物理逃逸速度$v_{\rm esc}=[2(1-\Gamma_e)GM/R]^{1/2}$，在$T\gtrsim2.1\times10^4$ K原题取$v_\infty=v_{\rm esc}/2.6$。
(a)忽略辐射压，画$L=10^5,3\times10^5,10^6,2\times10^6\Ls$下$\log\dot M$对$M$；
(b)取$\Gamma_e=\kappa L/(4\pi cGM)$、$\kappa=0.3\ \mathrm{cm^2g^{-1}}$重画，说明影响。
\end{exercise}
\begin{solution}
记太阳单位$m,l,r$，逃逸速度$617.7\sqrt{(1-\Gamma_e)m/r}$ km/s。乘题式的$\sqrt r$后半径消掉，所以
\[
 \dot M=\frac{10^{-1.37}(l/10^6)^{2.07}}
 {(617.7/2.6)\sqrt m\sqrt{1-\Gamma_e}}\ \Ms/\mathrm{yr}.
\]
原题使用除2.6的约定，不因其他资料常见另一标定而无说明改成乘2.6。无辐射时$\dot M\propto m^{-1/2}$；同质量加入辐射使$\dot M$乘$(1-\Gamma_e)^{-1/2}$。实际四个光度的两组曲线：
\begin{center}\begin{tikzpicture}\begin{axis}[width=.76\textwidth,height=5.4cm,xmode=log,ymode=log,xlabel={$M/\Ms$},ylabel={$\dot M$（$\Ms$/yr）},legend style={font=\footnotesize},legend pos=north east,unbounded coords=jump]
\addplot[black,dashed] table[x=M_solar,y=Mdot_no_Gamma,col sep=comma]{qa/stellar-wind-L100000.csv};\addlegendentry{$10^5$，无辐射}
\addplot[black] table[x=M_solar,y=Mdot_Gamma,col sep=comma]{qa/stellar-wind-L100000.csv};\addlegendentry{$10^5$，含辐射}
\addplot[gray,dashed] table[x=M_solar,y=Mdot_no_Gamma,col sep=comma]{qa/stellar-wind-L300000.csv};\addlegendentry{$3\times10^5$，无辐射}
\addplot[gray] table[x=M_solar,y=Mdot_Gamma,col sep=comma]{qa/stellar-wind-L300000.csv};\addlegendentry{$3\times10^5$，含辐射}
\addplot[structurecolor,dashed] table[x=M_solar,y=Mdot_no_Gamma,col sep=comma]{qa/stellar-wind-L1000000.csv};\addlegendentry{$10^6$，无辐射}
\addplot[structurecolor] table[x=M_solar,y=Mdot_Gamma,col sep=comma]{qa/stellar-wind-L1000000.csv};\addlegendentry{$10^6$，含辐射}
\addplot[main,dotted] table[x=M_solar,y=Mdot_no_Gamma,col sep=comma]{qa/stellar-wind-L2000000.csv};\addlegendentry{$2\times10^6$，无辐射}
\addplot[main,thick] table[x=M_solar,y=Mdot_Gamma,col sep=comma]{qa/stellar-wind-L2000000.csv};\addlegendentry{$2\times10^6$，含辐射}
\end{axis}\end{tikzpicture}\end{center}
$\Gamma_e<1$才有实逃逸速度，含辐射曲线只画满足该域的点。$\Gamma_e\to1^-$时经验式形式发散，表示假设失效，不是恒星可维持无限质量流；$\Gamma_e\geq1$不画成复数速度曲线。
\end{solution}

\src{8.901 PS6.5，冷亮星的Reimers风，(a)--(d)。}
\begin{exercise}
题给$\dot M=4\times10^{-13}\eta L/(gR)\ \Ms$/yr，$L,R,g$均为太阳单位，$\eta$约为1。
(a)估$M=1\Ms,L=7000\Ls,T=3000$ K的AGB星失质量率；
(b)错误地假定$L,R,\eta$恒定，求$M(t)$；
(c)用$L=7000,R=310,M_0=1,\eta=1$画质量对$10^5$ yr单位时间；
(d)多久减至$0.6\Ms$碳氧核心？
\end{exercise}
\begin{solution}
太阳单位下$g=m/r^2$，给定流失量的正值为$4\times10^{-13}\eta lr/m$，恒星质量导数为负数。(a)光度和温度给$r=\sqrt{7000}(5772/3000)^2=309.7$，初始流失率约$8.67\times10^{-7}\Ms$/yr。

(b) $t$用yr、$m=M/\Ms$，有$m\dot m=-4\times10^{-13}\eta lr$，故
\[
 m(t)=\sqrt{m_0^2-8\times10^{-13}\eta lr\,t}.
\]
不能保持初始$g$不变再算线性失质量，因为$g$也随质量变。
(c)--(d) 按题指定$r=310$，到$m_f=0.6$需
\[
 t_f=\frac{1-0.6^2}{8\times10^{-13}\times7000\times310}
 =3.687\times10^5\ \mathrm{yr}.
\]
\begin{center}\begin{tikzpicture}\begin{axis}[width=.6\textwidth,height=4.3cm,xlabel={$t$（$10^5$ yr）},ylabel={$M/\Ms$},ymin=.55,ymax=1.02]
\addplot[structurecolor] table[x=t_1e5yr,y=M_solar,col sep=comma]{qa/stellar-reimers.csv};
\end{axis}\end{tikzpicture}\end{center}
在给定核质量停止。更晚平方根可失去实数解，是固定光度半径模型被过度外推，不能解释为真实核心质量为负或恒星突然全部消失。
\end{solution}
